\section{Source audit, exact verification, and integration decisions}
\label{sec:audit}

\subsection{What changes in the source status}

The article supplies advances at specific research interfaces. It does
not replace the source's definitions or silently promote other open
claims. The following status changes are proposed for integration.

\begingroup
\small
\begin{longtable}{@{}p{0.27\textwidth}p{0.65\textwidth}@{}}
\toprule
Source target & Proposed status after review of this article\\
\midrule
\endhead
Quartic report, R4: fifth- and higher-power cancellation &
The universal recurrence and anchored generator are proved in
Theorems~\ref{tp:thm:recurrence} and \ref{tp:thm:generator}. Exact
eight- and eleven-coordinate rows are supplied for powers five and six.
Theorem~\ref{et:main} treats every repeated even order. The minimal
general coordinate count remains a separate question.\\[5pt]
Quartic report, R5: resonant parity &
The complete Laurent normalization and all coefficients are proved in
Theorem~\ref{tp:thm:Laurent}. The local spectral Taylor series has exact
radius two. A uniform finite arithmetic reduction as the resonance
index varies is not proved.\\[5pt]
Quartic report, R1: smaller order ten &
The complete all-$b$ statement for $b\ge6$ is proved in
Theorem~\ref{mt10:main}; the new missing certificate is $b=962$.
The smaller-order-twelve and unrestricted unequal-order questions remain
open here.\\[5pt]
Harmonic-gap continuation, Q4 and Q5 &
The first two-marker target and all identical marked clusters are
proved by \eqref{gm:gamma-square} and \eqref{gm:identical}.
Corollary~\ref{gm:marked} treats the fully symmetric unequal-marker
component. General individual ordered unequal-marker reductions remain
open.\\[5pt]
Higher-defect report, face and odd-weight questions &
Theorems~\ref{sd:thm:master}--\ref{sd:thm:allorders} give the ordinary
face decomposition and its full rational zero-jet term. The periodic
extension is Corollary~\ref{sd:cor:periodic}. Special-function evaluation
of general higher face integrals remains open.\\
\bottomrule
\end{longtable}
\endgroup

The general all-power generator is not the already known finite-span
theorem for powers of the ordinary harmonic number $H_n$. Here the
objects are powers of the decaying polygamma tail $\zeta(2,n+1)$.
Likewise the face theorem is not a new claim that a Dirichlet mixture
exists: it is the explicit decomposition and integrability of its
boundary-subtracted pieces. The relative Hurwitz interpolation itself
is inherited from published work. These separations prevent duplicate
claims during integration.

\subsection{Corrections and normalization details}

No new false theorem was found in the targeted rank and harmonic-gap
material reviewed for this continuation. In particular, the source's
single-marker theorem correctly left the several-marker extension open.
The lower-index-zero binomial omission in the reconstruction matrices
was already explicit in the canonical fifth-order report and in the
incoming extensions. It must be retained as an existing convention,
not reported as a newly discovered erratum.

The following details are essential when the new results are inserted.
They are precise safeguards for the supplied proof chain.

\begin{enumerate}
\item In the centered finite-cutoff argument use $n=0,\ldots,N-1$.
The even-power subtraction sums then have zero constant coefficient.
Changing to a cutoff ending at $N$ requires its endpoint polynomial.
The missing $n=0$ endpoint also produces the exceptional
$-d\delta_{k,2}\zeta(2)^{d-1}/2$ in
\eqref{tp:eq:recurrence}.

\item The tail expansions are asymptotic. Only a finite subtraction,
with a normally convergent remainder, is used at any specified
spectral point. A fixed infinite asymptotic series is not summed as
an ordinary power series in $1/n$.

\item The words ``eight coordinates,'' ``eleven coordinates,'' and
``rank $2b$'' concern the indicated rational coefficient maps. They do
not assert that the corresponding real or complex constants are
linearly independent over $\mathbb Q$. The source's existing formal
rank convention is preserved.

\item In the multiple-mark formula differentiate independent parameters
before specialization. A total derivative along the diagonal is a
different expression. For unequal marks, division of the symmetric
sum by $k!$ is not valid. For identical marked pairs it is valid.

\item The Gamma expressions use differences of analytic log-Gamma
branches on the right half-plane. An independently chosen logarithm
of a quotient need not have the same branch. The Stieltjes primitive
vanishes at its stated lower endpoint.

\item In the spectral decomposition the zero-jet correction
$\mathcal L_W$ is only the contribution of $F_0(B)$. Lower derivatives
$F_1,F_2,\ldots$ can still contain the Stieltjes and residue terms
present in the existing one-mass formulas. Removing $\mathcal L_W$
does not remove all such terms. The pair term $\zeta(2)$ in
\eqref{sd:eq:fourth} is required by the Gamma-ratio normalization.

\item A simplex face term is not an irreducible connected cumulant.
It can be nonzero when only two of several masses coincide. Formula
\eqref{sd:eq:quadraticface} gives its precise quadratic behavior and
prevents an incorrect stronger collision claim.
\end{enumerate}

An independent audit strengthened the proof of the polylogarithm
generator by replacing generic endpoint-regularization language with
an explicit convergent-MZV formula for the sole leading-one endpoint.
The Laurent derivative estimate was also written with its required
logarithmic factor. These are corrections made to the present draft
during review, not errors attributed to the baseline.

\subsection{What the reproducible checks establish}

The package separates proof certificates from numerical diagnostics.
The essential exact checks are as follows.

\begin{itemize}
\item A finite Faulhaber calculation, independent of the new Bernoulli
recurrence, agrees with it in 220 rational checks across powers one
through eight, including nonpositive anchors and boundary parity.
It retains MZVs as formal words. A separate finite-cutoff replay proves
the reduction of \eqref{tp:eq:main5} using only four classical MZV
relations. Rational coefficient extraction reproduces the full
eight-by-eight and eleven-by-eleven determinants and elimination rows.
For the repeated even-order extension, an independent standard-library
program checks 16 cubic cutoff polynomial identities and 192 finite
tail-difference versus strict-composition anchor identities.

\item The order-ten replay reconstructs the inherited determinant
polynomial from 120 signed permutation products. It checks all
$10{,}403$ integer monomials, the Bernoulli modular substitution,
and the small base matrices. A separate program computes all 25
coefficient norms used in the eventual argument and verifies the
integer factorial comparison. The new exceptional certificate uses
the displayed five-by-five matrix and independently computed Bernoulli
residues. No numerical determinant decides nonvanishing.

\item The marked-gap replay checks the formal insertion and partition
algebra through five independent labels. At five labels there are 52
set partitions and 120 surviving permutations; all unwanted endpoint
factors and shorter words cancel. The analytic insertion theorem is
proved in the text; the program checks its finite combinatorial
consequences.
\end{itemize}

Independent numerical representations then test the conclusions.
Three ordinary fifth- and sixth-power identities agree to more than
90 decimal digits using finite direct sums and centered Hurwitz tail
acceleration at two different cutoffs. The noninteger-endpoint marked
gap and its actual zero-marker derivative are reconstructed from
ordinary nested sums and compared with the decorated formulas; the
latter residual at cutoff 160 is approximately $3.73\times10^{-31}$.
Finite-interval Gamma and Stieltjes specializations agree to the stated
working precision. For the three-mass spectral third derivative, a
direct spectral evaluation agrees with independent Gamma face
quadrature to approximately $8.2\times10^{-32}$ at the highest
reported quadrature order.

These errors are observed cross-representation residuals. They are not
interval enclosures, and common truncation errors can survive a
comparison. The theorems rest on the analytic proofs and finite exact
identities. The package records the working precision, cutoff, method,
and exact-versus-numerical status of the checks. Its scripts use local
files only and need no repository checkout or network access to run.

The archive includes a standalone \LaTeX{} source, the compiled PDF,
modular section sources, Python verification programs, successful output
records, a claim-status ledger, source hashes, integration notes, and a
manifest verifier. The standalone source is generated from the modular
article, so the readable PDF and reusable section fragments share the
same mathematics. The source repository is left unchanged.

\section{Further research questions and exact next targets}
\label{sec:research}

The most useful next questions ask for identities, finite reductions,
or proofs of specified families. The following targets arise directly
from the present theorems.

\begin{enumerate}[label=\textbf{R\arabic*.}]
\item \textbf{Determine the smallest formal coordinate system at every
tail power.}
The universal generator supplies a finite span and exact ranks eight
and eleven for the chosen fifth- and sixth-power systems. Construct
canonical reduced coordinates at power $d$, state the admitted MZV
relations, and prove the rank of the resulting moment coefficient map.
One should distinguish the rank of that formal map from the unknown
arithmetic dimension after evaluating its coordinates.

\item \textbf{Classify compact polynomial annihilators.}
Equations \eqref{tp:eq:main5} and \eqref{tp:eq:main6} cancel all new
high-weight anchors in selected moment combinations. Find a structural
description of the polynomial weights producing a prescribed lower
coordinate span. A useful theorem would give the least degree and a
closed basis of such weights, with their exact asymptotic subtraction
polynomials. This could explain why the relatively small coefficients
in \eqref{intro:quintic} exist.

\item \textbf{Reduce the new anchored generators to explicit
polylogarithmic forms.}
The generator theorem is a finite algorithm with a depth upper bound.
Write a compact explicit expression for $H_5$ or $H_6$, then determine
which identities among allowed arguments and products reduce it.
Theorem statements should specify the function class and branches.
A proof of minimal depth would require an appropriate functional
invariant and is not supplied by the finite coordinate determinants.

\item \textbf{Find arithmetic recurrences for the resonant Laurent
family.}
Theorem~\ref{tp:thm:Laurent} fixes every coefficient and its spectral
germ. Determine whether the dependence on $r$ can be described by an
anchored recurrence with a controlled set of new logarithmic constants.
The first concrete case is the quartic constant term
$\Lambda_{4,r,0}$. Any proposed recurrence must reproduce the
$h_{4,r}\gamma_0(1/2)$ contribution exactly.

\item \textbf{Complete smaller order twelve, then remove the fixed-order
restriction.}
The order-ten proof shows how a modular certificate and a proved
eventual range can complement each other. Carry out the six-dimensional
contraction for smaller order twelve, account for every exceptional
residue below a sufficient threshold, and provide finite exact
certificates. A uniform determinant identity preventing cancellation
for all unequal even orders would settle the broader optimal-prefix
conjecture. The full gap-six family is still an intermediate target.

\item \textbf{Find closed reconstruction rows.}
The ordinary identities \eqref{mt10:ordinary} are explicit finite
rational algorithms involving an inverse matrix. Seek a fixed-size
kernel or product formula for its rows, including all endpoint
constants. This is a stronger goal than nonsingularity and may yield
short identities at large orders without constructing a large inverse.

\item \textbf{Resolve the antisymmetric two-marker component.}
For unequal marked pairs, the partition theorem gives the symmetric
combination. Determine an independent expression for
$\mathcal H(s,u,v,t)-\mathcal H(s,v,u,t)$ and its mixed derivatives.
A proposed reduction should retain the given relative quotient at
noninteger endpoints and prove an ordinary convergence domain before
continuation. This is the precise part of the several-marker question
not supplied by a classical symmetric-sum law.

\item \textbf{Combine marked clusters with nonpositive deletions.}
Unmarked deletions can shift distinct marked centers differently.
Classify exactly when such operations preserve the identical-cluster
or symmetric span, and identify the additional ordered coordinates
when they do not. The natural first test is one nonpositive integer
letter adjacent to a two-marker cluster. Independent differentiation
must precede every diagonal specialization.

\item \textbf{Evaluate special Gamma face integrals.}
For rational shifts and selected algebraic masses, seek identities
expressing $I_S(F_1)$ in a specified polylogarithmic, multiple-Gamma,
or zeta-derivative class. The quadratic identity
\eqref{sd:eq:quadraticface} gives a normalization check for every face
size. A finite exact integral is already available at arbitrary masses;
a special-function reduction would be an additional theorem.

\item \textbf{Give higher collision coefficients face by face.}
The quadratic coefficient in \eqref{sd:eq:collisionleading} is known
for every $k$. Derive closed symmetric-polynomial expressions for
$I_S(B^h)$ for general $h$, and for the logarithmically weighted
integrals in the higher derivative formula. Determine a generating
function that remains regular when any subset of masses coalesces.
This asks for explicit identities rather than only estimates.

\item \textbf{Allow independent spectral orders and arithmetic twists.}
Replace the common $s$ by independent $s_i$ in the product spectrum.
Derive the face decomposition for the resulting nonsymmetric Dirichlet
parameters, then compute prescribed transverse derivatives at their
joint origin. For periodic weights, project the resulting Gamma and
Stieltjes formulas onto Dirichlet characters while retaining the full
logarithmic normalization in the period. The mean-zero cancellation
provides an exact test at every polynomial weight.

\item \textbf{Return to the Gaussian period targets through an explicit
map.}
The new higher-depth polylogarithmic generators may supply functional
identities outside previously explored relation spaces. Before claiming
progress on $S_6$ or revised $S_8$, construct a precise map from one
such identity to the canonical residual, including its lower-weight
products and branch constants. A formal obstruction in a chosen
relation space and a high-precision small residual answer different
questions; neither alone proves or disproves the period identity.

\item \textbf{Formalize the finite-to-analytic interfaces.}
The Bernoulli boundary polynomial, finite summation identities,
modular determinant certificate, set-partition inversion, and face
M\"obius inversion are suitable finite formalization targets.
Separately formalize normal convergence of the subtracted spectral
families and agreement on their defining half-planes. Validated
interval quadrature would strengthen the computational diagnostics,
but remains distinct from formal proof of the exact identities.
\end{enumerate}
